\chapter{Vektor dan Garis pada Bidang}\label{ch:g11:vect}

\omterm{def:g11:vect:vector}{Vektor} menyandikan perpindahan; garis adalah lintasannya. Perkakas baru yang
penting pada bab ini adalah \emph{\omterm{def:g11:vect:det}{determinan}}, sebuah bilangan tunggal yang
memutuskan apakah dua \omterm{def:g11:vect:collinear}{vektor kolinear} --- lalu, dari sana, menghasilkan
\omterm{def:g10:algebra:equation}{persamaan} garis serta menuntaskan pertanyaan kesegarisan dan kesejajaran lewat
perhitungan semata. Gagasan yang sama diperluas ke ruang pada \cref{ch:g12:space}.

\section{Vektor: ringkasan}

\begin{definition}[Vektor, koordinat]\label{def:g11:vect:vector}
Sebuah \emph{vektor}\index{vektor} $\vec u$ menyatakan perpindahan: semua
anak panah dengan arah, jurusan dan panjang yang sama menyatakan vektor yang
sama. Pada sebuah \omterm{def:g10:coordgeom:system}{sistem koordinat}, $\vec u\,(x, y)$ mencatat perpindahannya
($x$ mendatar, $y$ tegak); untuk titik $A(x_A, y_A)$ dan
$B(x_B, y_B)$,
\[
\vect{AB}\,\bigl(x_B - x_A,\ y_B - y_A\bigr).
\]
Vektor dijumlahkan \omterm{def:g10:coordgeom:system}{koordinat} demi \omterm{def:g10:coordgeom:system}{koordinat}, dan $\lambda\vec u$ berkoordinat
$(\lambda x, \lambda y)$.
\end{definition}

\begin{proposition}[Titik tengah]\label{prop:g11:vect:midpoint}
\omterm{prop:g10:coordgeom:midpoint}{Titik tengah} $M$ ruas garis $[AB]$ berkoordinat
$\left(\frac{x_A + x_B}{2},\ \frac{y_A + y_B}{2}\right)$.
\end{proposition}

\begin{proof}
$M$ dicirikan oleh $\vect{AM} = \frac12\vect{AB}$: \omterm{def:g10:coordgeom:system}{koordinat} demi \omterm{def:g10:coordgeom:system}{koordinat},
$x_M - x_A = \frac12(x_B - x_A)$, jadi $x_M = \frac{x_A + x_B}{2}$, dan
serupa untuk $y_M$.
\end{proof}

\section{Kekolinearan dan determinan}

\begin{definition}[Vektor kolinear]\label{def:g11:vect:collinear}
Dua \omterm{def:g11:vect:vector}{vektor} disebut \emph{kolinear}\index{vektor kolinear} bila yang satu
merupakan kelipatan skalar yang lain: $\vec v = \lambda \vec u$ untuk suatu
$\lambda \in \R$ (atau $\vec u = \vec 0$). Secara geometris: keduanya mempunyai
arah yang sama, atau salah satunya nol.
\end{definition}

\begin{definition}[Determinan]\label{def:g11:vect:det}
\emph{Determinan}\index{determinan} $\vec u\,(x, y)$ dan
$\vec v\,(x', y')$ adalah
\[
\det(\vec u, \vec v) =
\begin{vmatrix} x & x' \\ y & y' \end{vmatrix}
= x y' - x' y .
\]
\end{definition}

\begin{theorem}[Kriteria kekolinearan]\label{thm:g11:vect:collinearity}
Dua \omterm{def:g11:vect:vector}{vektor} $\vec u$ dan $\vec v$ \omterm{def:g11:vect:collinear}{kolinear} jika dan hanya jika
$\det(\vec u, \vec v) = 0$.
\end{theorem}

\begin{proof}
($\Rightarrow$) Jika $\vec v = \lambda\vec u$, maka $x' = \lambda x$ dan
$y' = \lambda y$, sehingga
$\det(\vec u, \vec v) = x(\lambda y) - (\lambda x)y = 0$; bila
$\vec u = \vec 0$, \omterm{def:g11:vect:det}{determinannya} jelas $0$.

($\Leftarrow$) Andaikan $xy' - x'y = 0$ dan $\vec u \neq \vec 0$, katakan
$x \neq 0$ (kasus $y \neq 0$ setangkup). Ambil
$\lambda = \frac{x'}{x}$. Maka $x' = \lambda x$ menurut susunannya, dan
hubungan $xy' = x'y = \lambda x y$ memberi, setelah dibagi $x \neq 0$,
$y' = \lambda y$. Jadi $\vec v = \lambda \vec u$.
\end{proof}

\begin{example}\label{ex:g11:vect:collinear}
$\vec u\,(3, -2)$ dan $\vec v\,(-6, 4)$:
$\det = 3 \times 4 - (-6) \times (-2) = 12 - 12 = 0$, jadi \omterm{def:g11:vect:collinear}{kolinear} (memang
$\vec v = -2\vec u$). Untuk $\vec u\,(3, -2)$ dan $\vec w\,(2, 1)$:
$\det = 3 + 4 = 7 \neq 0$, jadi tidak \omterm{def:g11:vect:collinear}{kolinear}.
\end{example}

\begin{method}[Kesegarisan tiga titik]\label{met:g11:vect:alignment}
Titik $A$, $B$, $C$ segaris jika dan hanya jika $\vect{AB}$ dan $\vect{AC}$
\omterm{def:g11:vect:collinear}{kolinear}: hitunglah kedua \omterm{def:g11:vect:vector}{vektornya} lalu periksalah
$\det\bigl(\vect{AB}, \vect{AC}\bigr) = 0$.
\end{method}

\section{Persamaan Kartesius sebuah garis}

\begin{definition}[Vektor arah]\label{def:g11:vect:direction}
Sebuah \emph{vektor arah}\index{vektor arah} garis $d$ adalah sebarang
\omterm{def:g11:vect:vector}{vektor} taknol $\vec u$ yang membuat $\vect{AB}$ \omterm{def:g11:vect:collinear}{kolinear} dengan $\vec u$
untuk semua titik $A, B$ pada $d$.
\end{definition}

\begin{theorem}[Persamaan Kartesius]\label{thm:g11:vect:cartesian}
Sebuah garis dengan \omterm{def:g11:vect:direction}{vektor arah} $\vec u\,(-b, a)$ mempunyai \omterm{def:g10:algebra:equation}{persamaan}
berbentuk
\[
ax + by + c = 0
\qquad (a, b) \neq (0, 0).
\]
Sebaliknya, setiap \omterm{def:g10:algebra:equation}{persamaan} semacam itu memerikan garis dengan \omterm{def:g11:vect:direction}{vektor arah}
$\vec u\,(-b, a)$.
\end{theorem}

\begin{proof}
Misalkan $d$ melalui $A(x_A, y_A)$ dengan arah $\vec u\,(-b, a)$. Sebuah
titik $M(x, y)$ terletak pada $d$ tepat ketika $\vect{AM}$ \omterm{def:g11:vect:collinear}{kolinear} dengan
$\vec u$, yaitu (\cref{thm:g11:vect:collinearity})
\[
0 = \det\bigl(\vect{AM}, \vec u\bigr)
= (x - x_A)\,a - (-b)(y - y_A)
= ax + by - (a x_A + b y_A),
\]
sebuah \omterm{def:g10:algebra:equation}{persamaan} berbentuk seperti yang disebutkan, dengan $c = -(a x_A + b y_A)$.
Sebaliknya, penyelesaian $ax + by + c = 0$ dengan, katakanlah, $b \neq 0$ adalah
titik $\left(x, -\frac{a x + c}{b}\right)$: mengurangkan dua di antaranya
memperlihatkan bahwa setiap tali busurnya \omterm{def:g11:vect:collinear}{kolinear} dengan $(-b, a)$, dan selalu
ada satu penyelesaian --- jadi itulah garis melaluinya dengan arah $(-b, a)$.
\end{proof}

\begin{method}[Garis melalui dua titik]\label{met:g11:vect:twopoints}
Untuk mencari \omterm{def:g10:algebra:equation}{persamaan} garis $(AB)$: hitunglah \omterm{def:g11:vect:direction}{vektor arahnya}
$\vect{AB}\,(x_B - x_A, y_B - y_A)$; padankan dengan $(-b, a)$, yaitu ambil
$a = y_B - y_A$ dan $b = -(x_B - x_A)$; lalu tentukan $c$ dengan memasukkan
\omterm{def:g10:coordgeom:system}{koordinat} $A$.
\end{method}

\begin{example}\label{ex:g11:vect:twopoints}
Garis melalui $A(1, 2)$ dan $B(4, 3)$: $\vect{AB}\,(3, 1)$, sehingga
$a = 1$, $b = -3$, dan \omterm{def:g10:algebra:equation}{persamaannya} $x - 3y + c = 0$. Memasukkan $A$:
$1 - 6 + c = 0$, jadi $c = 5$. \omterm{def:g10:algebra:equation}{Persamaannya}: $x - 3y + 5 = 0$. Periksa dengan
$B$: $4 - 9 + 5 = 0$.
\end{example}

\begin{remark}
Ketika $b \neq 0$, \omterm{def:g10:algebra:equation}{persamaannya} dapat diselesaikan untuk $y$: $y = mx + p$
dengan \omterm{def:g10:reffunc:affine}{gradien} $m = -\frac ab$. Bentuk Kartesius $ax + by + c = 0$ lebih
umum: bentuk itu juga meliputi garis tegak ($b = 0$), yang tidak bergradien.
\end{remark}

\begin{omfigure}
\begin{tikzpicture}
\begin{axis}[omaxis,
  width=8.6cm, height=5.8cm,
  xmin=-1.6, xmax=5.6, ymin=-0.9, ymax=4.4,
  xtick={1,4}, ytick={2,3}]
  \addplot[omDef, thick, domain=-1.4:5.4] {(x+5)/3};
  \fill (axis cs:1,2) circle (1.5pt)
    node[above left, font=\small] {$A$};
  \fill (axis cs:4,3) circle (1.5pt)
    node[above left, font=\small] {$B$};
  \draw[-Stealth, omThm, thick] (axis cs:1,2) -- (axis cs:2.5,2.5)
    node[midway, below right, font=\small] {$\vec u$};
\end{axis}
\end{tikzpicture}

{\small Garis $x - 3y + 5 = 0$ pada \cref{ex:g11:vect:twopoints}, beserta
\omterm{def:g11:vect:direction}{vektor arahnya} $\vec u = \vect{AB}$ (merah, di sini terskalakan $\frac12$),
dan kedua titik yang dipakai untuk menyusun \omterm{def:g10:algebra:equation}{persamaannya}.}
\end{omfigure}

\section{Perian parametrik}

\begin{definition}[Perian parametrik]\label{def:g11:vect:parametric}
Garis melalui $A(x_A, y_A)$ dengan arah $\vec u\,(\alpha, \beta)$ adalah
himpunan titik
\[
M\bigl(x_A + t\alpha,\ y_A + t\beta\bigr), \qquad t \in \R .
\]
Bilangan $t$ disebut \emph{parameternya}: setiap nilai $t$ menghasilkan satu
titik pada garis itu, seolah-olah menyusurinya dengan kecepatan tetap $\vec u$.
\end{definition}

\begin{example}
Garis pada \cref{ex:g11:vect:twopoints} juga sama dengan
$\{(1 + 3t,\ 2 + t) : t \in \R\}$: $t = 0$ memberi $A$, $t = 1$ memberi $B$.
Melenyapkan $t$ ($t = y - 2$, sehingga $x = 1 + 3(y-2)$) mengembalikan
$x - 3y + 5 = 0$.
\end{example}

\section{Kedudukan relatif dua garis}

\begin{proposition}[Sejajar atau berpotongan]\label{prop:g11:vect:positions}
Dua garis dengan \omterm{def:g11:vect:direction}{vektor arah} $\vec u$ dan $\vec v$ sejajar jika dan
hanya jika $\det(\vec u, \vec v) = 0$. Dalam bentuk \omterm{def:g10:algebra:equation}{persamaan}:
$d : ax + by + c = 0$ dan $d' : a'x + b'y + c' = 0$ sejajar jika dan hanya jika
$ab' - a'b = 0$. Garis yang tidak sejajar bertemu di tepat satu titik.
\end{proposition}

\begin{proof}
Kesejajaran berarti arahnya \omterm{def:g11:vect:collinear}{kolinear}, dan itulah
\cref{thm:g11:vect:collinearity}; dengan $\vec u\,(-b, a)$ dan
$\vec v\,(-b', a')$, \omterm{def:g11:vect:det}{determinannya} adalah $(-b)a' - (-b')a = ab' - a'b$.
Jika kedua garisnya tidak sejajar, menyelesaikan kedua \omterm{def:g10:algebra:equation}{persamaannya} serentak
(lewat substitusi atau kombinasi) menghasilkan penyelesaian tunggal: sistemnya
berperilaku seperti sistem $2 \times 2$ berdeterminan taknol.
\end{proof}

\begin{method}[Perpotongan dua garis]\label{met:g11:vect:intersection}
Selesaikan sistem kedua \omterm{def:g10:algebra:equation}{persamaannya}: pisahkan satu variabel pada salah satu
\omterm{def:g10:algebra:equation}{persamaan} lalu substitusikan ke \omterm{def:g10:algebra:equation}{persamaan} lainnya, atau padukan kedua
\omterm{def:g10:algebra:equation}{persamaannya} untuk melenyapkan satu variabel. Selalu periksa titik hasilnya
pada \emph{kedua} \omterm{def:g10:algebra:equation}{persamaan}.
\end{method}

\begin{example}\label{ex:g11:vect:intersection}
Potongkan $d: x - 3y + 5 = 0$ dan $d': 2x + y - 4 = 0$. Dari $d'$:
$y = 4 - 2x$. Setelah disubstitusikan ke $d$:
\[
x - 3(4 - 2x) + 5 = 0
\iff x - 12 + 6x + 5 = 0
\iff 7x = 7
\iff x = 1,
\]
lalu $y = 2$. Kedua garisnya bertemu di $(1, 2)$. Periksa pada $d$: $1 - 6 + 5 = 0$.
\end{example}

\section{Latihan}

\begin{exercise}[$\star$]\label{exo:g11:vect:1}
Diberikan $A(2, -1)$, $B(5, 3)$ dan $C(-1, 1)$, hitunglah \omterm{def:g10:coordgeom:system}{koordinat}
$\vect{AB}$, $\vect{AC}$, $\vect{AB} + \vect{AC}$ dan $2\vect{AB} -
3\vect{AC}$, serta \omterm{prop:g10:coordgeom:midpoint}{titik tengah} $[BC]$.
\end{exercise}

\begin{exercise}[$\star$]\label{exo:g11:vect:2}
Apakah \omterm{def:g11:vect:vector}{vektor} berikut \omterm{def:g11:vect:collinear}{kolinear}?
\[
\vec u\,(4, -6) \text{ dan } \vec v\,(-2, 3); \qquad
\vec u\,(2, 5) \text{ dan } \vec v\,(3, 7); \qquad
\vec u\,(0, 3) \text{ dan } \vec v\,(0, -8).
\]
\end{exercise}

\begin{exercise}[$\star$]\label{exo:g11:vect:3}
Carilah \omterm{def:g10:algebra:equation}{persamaan} Kartesius garis melalui $A(2, 1)$ dengan \omterm{def:g11:vect:direction}{vektor
arah} $\vec u\,(3, -1)$, lalu garis melalui $B(0, 4)$ dan
$C(2, 0)$.
\end{exercise}

\begin{exercise}[$\star$]\label{exo:g11:vect:4}
Untuk garis $d: 3x - 2y + 6 = 0$, berikan sebuah \omterm{def:g11:vect:direction}{vektor arah},
titik potongnya dengan kedua sumbu \omterm{def:g10:coordgeom:system}{koordinat}, lalu putuskan apakah titik
$P(2, 6)$ dan $Q(1, 4)$ termasuk pada $d$.
\end{exercise}

\begin{exercise}[$\star\star$]\label{exo:g11:vect:5}
Apakah titik $A(1, 2)$, $B(4, 8)$, $C(-2, -4)$ segaris? Pertanyaan yang sama
untuk $D(0, 1)$, $E(2, 4)$, $F(5, 9)$.
\end{exercise}

\begin{exercise}[$\star\star$]\label{exo:g11:vect:6}
Tentukan titik potongnya, bila ada:
\[
d: 2x + y - 7 = 0 \ \text{ dan }\ d': x - y + 1 = 0;
\qquad
e: x - 2y + 3 = 0 \ \text{ dan }\ e': -2x + 4y + 1 = 0 .
\]
\end{exercise}

\begin{exercise}[$\star\star$]\label{exo:g11:vect:7}
Berikan perian parametrik garis $d: 2x - y + 3 = 0$, lalu sebuah
\omterm{def:g10:algebra:equation}{persamaan} Kartesius garis
$\bigl\{(1 + 2t,\ -3 + t) : t \in \R\bigr\}$.
\end{exercise}

\begin{exercise}[$\star\star$]\label{exo:g11:vect:8}
Carilah \omterm{def:g10:algebra:equation}{persamaan} garis melalui $P(1, -2)$ yang sejajar dengan
$d: 4x - 3y + 2 = 0$, lalu nilai $m$ yang membuat garis
$mx + 2y - 5 = 0$ sejajar dengan $d$.
\end{exercise}

\begin{exercise}[$\star\star$]\label{exo:g11:vect:9}
Misalkan $A(-1, 0)$, $B(3, 2)$ dan $C(1, -4)$. Carilah \omterm{def:g10:algebra:equation}{persamaan} Kartesius
\emph{garis berat} segitiga $ABC$ yang ditarik dari $C$ (yaitu garis yang
menghubungkan $C$ dengan \omterm{prop:g10:coordgeom:midpoint}{titik tengah} $[AB]$).
\end{exercise}

\begin{exercise}[$\star\star$]\label{exo:g11:vect:10}
Untuk nilai parameter $m$ yang mana \omterm{def:g11:vect:vector}{vektor}
$\vec u\,(m, 2)$ dan $\vec v\,(3, m - 1)$ \omterm{def:g11:vect:collinear}{kolinear}?
\end{exercise}

\begin{exercise}[$\star\star\star$]\label{exo:g11:vect:11}
Misalkan $ABCD$ sebuah jajargenjang, $I$ \omterm{prop:g10:coordgeom:midpoint}{titik tengah} $[AB]$, dan $J$
titik yang ditentukan oleh $\vect{DJ} = \frac23 \vect{DI}$. Bekerjalah pada
\omterm{def:g10:coordgeom:system}{sistem koordinat} dengan $A(0,0)$, $B(1,0)$, $D(0,1)$ (sehingga $C(1,1)$):
\begin{enumerate}
  \item Hitunglah \omterm{def:g10:coordgeom:system}{koordinat} $I$ dan $J$.
  \item Tunjukkan bahwa $A$, $J$ dan $C$ segaris. (Sebuah fakta yang manis:
        garis $(DI)$ memotong diagonal $(AC)$ pada sepertiga jalannya.)
\end{enumerate}
\end{exercise}

\section{Soal: Haluan bertabrakan dan koordinat yang cerdik}

\begin{problem}[{Soal akhir pekan --- lintasan yang berpotongan bukan
berarti bertabrakan: jarak terdekat di laut, dan teorema yang dibuktikan
dengan memilih kerangka yang tepat}]
\label{pb:g11:vect:1}
Haluan lurus dua kapal berpotongan --- haruskah kedua kapal itu
bertabrakan? Tidak,
kecuali keduanya mencapai titik potong itu \emph{pada waktu yang sama}:
lintasan adalah himpunan titik, gerak bersifat parametrik, dan mengacaukan
keduanya sudah menenggelamkan kapal sungguhan. Soal ini menjalankan sebuah
layanan lalu lintas kapal kecil pada garis parametrik
(\cref{def:g11:vect:parametric}), lalu beralih ke geometri murni dengan
kekuatan super \omterm{def:g10:coordgeom:system}{koordinat} yang lain: \emph{memilih} \omterm{def:g10:coordgeom:system}{koordinatnya} secara cerdik,
seperti pada \cref{exo:g11:vect:11}, agar teorema klasik menghitung dirinya
sendiri.

\medskip
\noindent\textbf{Bagian I --- Garis, dengan tiga cara.}
\begin{enumerate}
  \item Berikan perian parametrik garis melalui
        $A(1, 2)$ dengan \omterm{def:g11:vect:direction}{vektor arah} $\vec u(3, -1)$, lalu
        lenyapkan parameternya untuk memperoleh \omterm{def:g10:algebra:equation}{persamaan} Kartesius.
  \item Untuk garis $2x - 5y + 3 = 0$: bacalah sebuah \omterm{def:g11:vect:direction}{vektor
        arahnya} (\cref{thm:g11:vect:cartesian}), carilah satu titik padanya,
        lalu tulislah perian parametriknya.
  \item Hitunglah titik potong kedua garis pada
        pertanyaan 1 dan 2 (\cref{met:g11:vect:intersection}).
  \item Benarkan dengan \omterm{def:g11:vect:det}{determinannya}
        (\cref{thm:g11:vect:collinearity}) bahwa kedua garis itu
        memang harus berpotongan; lalu tunjukkan bahwa $2x - 5y + 3 = 0$ dan
        $-4x + 10y + 1 = 0$ sejajar secara tegas.
  \item Apakah $A(2, 1)$, $B(5, 3)$, $C(11, 7)$ segaris
        (\cref{met:g11:vect:alignment})?
\end{enumerate}

\medskip
\noindent\textbf{Bagian II --- Layanan lalu lintas kapal.}
Pada waktu $t$ (jam), kapal $P$ berada di $(2t,\ 3 + t)$ dan kapal $Q$
di $(10 - t,\ 2t - 1)$ (jaraknya dalam mil laut).
\begin{enumerate}[resume]
  \item Carilah \omterm{def:g10:algebra:equation}{persamaan} Kartesius kedua \emph{lintasannya},
        lalu hitunglah tempat kedua lintasan itu berpotongan.
  \item Pada waktu berapa masing-masing kapal melewati titik potongnya?
        Apakah kedua kapal itu bertabrakan? Nyatakan peringatan umumnya dalam
        satu kalimat: lintasan yang berpotongan melawan gerak yang berjumpa.
  \item Hitunglah kuadrat jarak $D^2(t)$ antara kedua kapal
        pada waktu $t$, lalu susutkan menjadi bentuk kuadrat dalam $t$. Pada
        waktu berapa kedua kapal itu terdekat, dan seberapa dekat
        keduanya (dengan teknologi titik puncak \cref{pb:g11:quad:1})?
  \item Peraturan menuntut pemisahan paling sedikit $1$ mil.
        Apakah aturan itu dilanggar? Jika ya, selesaikan $D^2(t) < 1$ lalu
        berikan lama pelanggarannya.
  \item Jelaskan mengapa, bagi \emph{sebarang} dua kapal pada
        haluan lurus dengan laju tetap, $D^2(t)$ selalu merupakan \omterm{def:g11:quad:quadratic}{fungsi
        kuadrat} dari $t$ --- sehingga ``titik pendekatan terdekat''
        selalu berupa perhitungan titik puncak.
  \item Pencegatan: sebuah kapal patroli berangkat dari titik asal pada
        $t = 0$ dengan kecepatan tetap $(a, b)$ dan harus menemui
        kapal $P$ tepat pada $t = 3$. Hitunglah
        $(a, b)$ yang diperlukan dan laju kapal patroli itu.
\end{enumerate}

\medskip
\noindent\textbf{Bagian III --- \omterm{def:g10:coordgeom:system}{Koordinat} yang cerdik.}
Bekerjalah pada kerangka \cref{exo:g11:vect:11}: jajargenjang
$ABCD$ menjadi $A(0,0)$, $B(1,0)$, $C(1,1)$, $D(0,1)$; misalkan $I$
\omterm{prop:g10:coordgeom:midpoint}{titik tengah} $[AB]$ dan $K$ \omterm{prop:g10:coordgeom:midpoint}{titik tengah} $[CD]$.
\begin{enumerate}[resume]
  \item Hitunglah \omterm{def:g10:algebra:equation}{persamaan} Kartesius garis $(DI)$.
  \item Potongkan $(DI)$ dengan diagonal $(AC)$: temukan kembali
        fakta \cref{exo:g11:vect:11} --- potongannya terletak tepat
        pada sepertiga jalan menaiki diagonalnya.
  \item Kini garis $(BK)$: carilah \omterm{def:g10:algebra:equation}{persamaannya}, potongkan dengan
        $(AC)$, lalu simpulkan teorema klasik lengkapnya:
        \emph{kedua ceva $(DI)$ dan $(BK)$ membagi tiga sama panjang
        diagonal $(AC)$}.
  \item Mengapa membuktikan teorema itu pada satu kerangka
        khusus ini sah? (Apa yang dipertahankan oleh pemilihan \omterm{def:g10:coordgeom:system}{koordinatnya},
        dan apa saja yang dilibatkan pernyataannya?) Jawablah
        dalam dua atau tiga kalimat.
  \item Satu pembagian tiga lagi sebagai latihan: misalkan $E$
        \omterm{prop:g10:coordgeom:midpoint}{titik tengah} $[BC]$. Di mana garis $(AE)$ memotong
        diagonal $(DB)$ yang \emph{satunya}?
\end{enumerate}

\medskip
\noindent\textbf{Bagian IV --- Teka-teki kedudukan.}
\begin{enumerate}[resume]
  \item Apakah ketiga garis $x + y = 4$, \ $2x - y = 2$, \
        $x - 2y = -2$ berpotongan di satu titik? (Potongkan dua di antaranya,
        lalu ujilah yang ketiga.)
  \item Untuk setiap \omterm{def:g10:numbers:sets}{bilangan real} $m$, misalkan $d_m$ garis
        $mx - y + 2 - 3m = 0$. Tunjukkan bahwa \emph{setiap} $d_m$
        melalui satu \omterm{pb:g10:functions:1}{titik tetap} yang sama. (Kelompokkan sukunya menurut
        $m$.)
  \item Pada keluarga $(d_m)$: anggota yang mana yang sejajar dengan
        $y = 2x + 7$? Yang mana yang melalui titik asal?
  \item Penutup --- ketiga busana sebuah garis: Kartesius
        (uji keanggotaan dengan satu substitusi), tereduksi (\omterm{def:g10:reffunc:affine}{gradien}
        dan \omterm{def:g10:functions:graph}{grafiknya} terbaca sekilas), parametrik (gerak, waktu,
        pencegatan); \omterm{def:g11:vect:det}{determinan} sebagai peramal kesejajaran;
        serta kerangka yang dipilih baik-baik sebagai pabrik teorema. Masing-masing
        satu kalimat.

\end{enumerate}
\end{problem}
